% ===================================================================
%  TABELLA nr. 7 — Il vitg durant l'enviern. Il bain rustical durant
%  la primavaira.
%  Traduction 2026 d'apres texto/io, controlee sur texto/fr, texto/ca
%  et texto/oc. Consigne : voir texto/rm/10-tabella-01.tex.
%
%  LA SECONDE SCENE MET EN FACE DEUX MANIERES DE NOMMER UNE BETE, ET
%  L'UNE DES DEUX EST CELLE DE L'IDO LUI-MEME.
%
%  Guignon DERIVE. Il pose un radical et le decline par le sexe et
%  par l'age : bov-ulo, bov-ino, bov-yuno ; muton-uli, muton-ini,
%  muton-yuni ; kapr-ini, kapro-yuni ; han-ulo, han-ino,
%  han-yuneti ; kaval-ino, kaval-yuno ; porkino, pork-yuno ;
%  anad-ini, anad-uli. Sept series, un seul procede.
%
%  LE ROMANCHE A UN MOT SEPARE POUR CHAQUE CASE :
%    « bov » le boeuf, « tor » le taureau, « vatga » la vache,
%      « vadè » le veau ;
%    « muntun » le belier, « nursa » la brebis, « agnè » l'agneau ;
%    « chaura » la chevre, « chavrin » le chevreau ;
%    « gial » le coq, « giaglina » la poule, « pulschin » le
%      poussin ;
%    « portg » le porc, « scrofa » la truie, « purtgrin » le
%      cochonnet ;
%    « chaval » le cheval, « juvla » la jument, « poltrin » le
%      poulain.
%
%  CE N'EST PAS UNE COQUETTERIE DE LEXIQUE. Sur un alpage la
%  difference entre une « nursa » et un « agnè » est la difference
%  entre deux travaux : l'une se tond et se trait, l'autre se porte.
%  Un mot compose oblige a lire deux morceaux avant de savoir quoi
%  faire ; un mot plein se comprend au cri.
%
%  ET C'EST LE TABLEAU 6 CONTINUE D'UN CRAN. Le 6 disait que la
%  richesse d'un vocabulaire suit L'OBJET — dix mots pour la
%  cheminee qu'on n'a pas, un compose pour le poele qu'on a. Le 7 dit
%  qu'elle suit LE TRAVAIL : le nombre de racines separees mesure la
%  distance a laquelle on se tient de la bete. Une langue construite
%  a Paris derive ; une langue parlee dans une etable nomme.
%
%  Le romanche du village, lui, se comporte comme partout ailleurs :
%  « faver », « chalzaner », « barbier », « charrader »,
%  « suedler », « mulinar », « pastrizer », « notar » — la meme
%  repartition qu'au tableau 5, la pierre latine et l'atelier venu
%  d'ailleurs. La difference est toute du cote de l'etable.
% ===================================================================

%%K t07-noto-1 noto
\VUnotes{143.2mm}{%
(*) Vesair per ils detagls dal past las tabellas 6 e 13.%
}

%%K t07-apar-1 apar
\VUsaut{4.0mm}
\VUcentre{13.2pt}{90}{SEGUNDA SERIA}
\VUsaut{3.0mm}
\VUfilet{16mm}
\VUsaut{5.6mm}
\VUtitre{15.0pt}{60}{TABELLA nr. 7}
\VUsaut{3.0mm}

%%K t07-tit-1 sub
\VUcentre{12.0pt}{0}{\textit{Emprima scena.}}
\VUsaut{3.0mm}

%%K t07-tit-2 sub
\VUpk{10.2pt}{40}{Il vitg durant l'enviern.}
\VUsaut{4.0mm}

%%K t07-c1-01-1 p
1. --- Durant las vacanzas da Nouv Onn hai jau accumpagnà mes bab a Vitgnov, in
pitschen vitg nua ch'el aveva da tractar intgins affars da commerzi.

%%K t07-c1-02-1 p
2. --- Nus avain duvrà la \VUgras{diligenza} \textsuperscript{(11)} che circulescha
tranter la citad e quel vitg; ma perquai ch'i faseva fitg fraid, n'è quel
viadi en in vehichel pauc confortabel betg stà fitg agreabel.

%%K t07-c1-03-1 p
3. --- Perquai avain nus sentì in ver plaschair cur ch'nus avain vis il
\VUgras{manader} fermar ses char sin la plazza, avant l'\VUgras{ustaria}
\textsuperscript{(2)} da l'\textit{Urs alv}. L'\VUgras{ustier} \textsuperscript{(4)} ans
spetgava sin il sulier da sia chasa, fimond cun quietezza sia \VUgras{pippa}.

%%K t07-c1-03-2 p
El ans ha envidà ad entrar, e jau na m'hai betg laschà rugar per m'installar avant
il fieu da sarments che crepitava en il fieughin. Alura hai jau fatg fitg gieu cun
l'asper \VUgras{vent da nord} che faseva crustgir la veglia \VUgras{ensaina}
\textsuperscript{(3)} e purtava davent en nairs vortischs \VUgras{il fim}
\textsuperscript{(1)} ch'ì ora dal chamin.

%%K t07-c1-04-1 p
4. --- I era l'ura dal gentar. Senza spetgar pli ditg avain nus maglià bain il past
simpel ma savurus ch'ins ans ha servì (*).

%%K t07-tit-3 sub
\VUpk{10.2pt}{40}{Intginas visitas.}
\VUsaut{3.0mm}

%%K t07-c1-05-1 p
5. --- Suenter il gentar hai jau accumpagnà mes bab, ch'è assicuratur, tar sia
clientella. L'emprim avain nus visità il \VUgras{faver} \textsuperscript{(5)}, che
abita sper l'ustaria. El era occupà da ferrar in chaval. Jau hai admirà la fermezza
da ses cumpogn, che tegneva solidamain sin ses scussal da corom l'\VUgras{unqua} dal
chaval, entant ch'il faver fixava il \VUgras{fier} \textsuperscript{(6)} cun ferms
colps da martè.

%%K t07-c1-06-1 p
6. --- Lura essan nus ids tar il \VUgras{chalzaner} \textsuperscript{(7)} dal vitg, il
vegl Jachen. Malgrà ch'el ha sco ensaina ina bellissima \VUgras{buttina}, sa
cuntentava el da metter novas soras ad in vegl pèr da chalzers pertgirads.

%%K t07-c1-06-2 p
Il vegl ans ha ditg ridend: «En la citad eri jau "chalzaner" e n'avevi nagins
clients. Qua sun jau "rapezzader da chalzers" e la lavur è bler.»

%%K t07-c1-07-1 p
7. --- El ans ha discurrì da ses vischin il \VUgras{barbier} \textsuperscript{(8)}, cun
il qual la clientella n'è betg cuntenta. Sias \VUgras{rasuras} èn adina sbrizzadas e
ses \VUgras{forschs} na tunden mai bain.

Ma perquai ch'el è l'unic barbier dal vitg, èn ins sforzads da sa laschar rasar e
tundrar da el. Dal rest è el in um avar e nungentil.

%%K t07-c1-08-1 p
8. --- Per fortuna n'assemeglian ils auters \VUgras{vischins} betg ad el. Il
\VUgras{charrader} \textsuperscript{(9)}, per exempel, è in um bunmanaivel, lavurus e
serviziaivel. I è in plaschair da laschar reparar in char da el. Sia lavur è adina
finida a temp.

%%K t07-c1-09-1 p
9. --- Er il vegl Jachen na s'è betg lamentà dal \VUgras{suedler}
\textsuperscript{(10)}, al qual el gida mintgatant a cuser ils \VUgras{guarnischs} cur
ch'la lavur pressa.

%%K t07-c1-09-2 p
Mes bab al ha dumandà da sia famiglia: «Tuts van bain. Mia abiadia è en la
\VUgras{scola} \textsuperscript{(12)}. Sia \VUgras{magistra} \textsuperscript{(13)} è
fitg cuntenta da ella, ella è ina buna \VUgras{scolara} \textsuperscript{(14)}. Ella
n'assemeglia betg a mes nauscha mat da figl; el pensa mo a giogar. Il
\VUgras{magister} \textsuperscript{(15)} na cuntanscha betg ch'el emprendia insatge.»

%%K t07-tit-4 sub
\VUsaut{3.0mm}
\VUpk{10.2pt}{40}{Tschintschadas dal vitg.}
\VUsaut{3.0mm}

%%K t07-c1-10-1 p
10. --- Il vegl ans ha lura raquintà las novitads dal vitg. Ins vegn a construir ina
nova scola, perquai che quella installada en la \VUgras{chasa communala} è daventada
memia pitschna. Dal rest è quai simpel, perquai ch'i basta da cumprar ina part da
l'iert dal \VUgras{mulinar}. Quai vegn a esser fitg profitaivel per il pauper um.
Pervi dal fraid na pon las \VUgras{muelas} dal \VUgras{malin} \textsuperscript{(16)}
betg pli maschinar il furment, perquai ch'la \VUgras{roda} \textsuperscript{(17)} è
vegnida fermada dal \VUgras{glatsch} \textsuperscript{(25)} da l'\VUgras{auatg}
\textsuperscript{(23)}.

%%K t07-c1-11-1 p
11. --- «Da construcziuns discurrend, avais Vus vis nossa nova \VUgras{baselgia}
\textsuperscript{(18)}? Ella è fitg pittoresca sut ses mantè da naiv. Ils
\VUgras{corvs} \textsuperscript{(19)}, che n'enconuschan betg pli lur vegl clutger,
sesan trists sin il grond arber dal \VUgras{santeri} \textsuperscript{(20)}.»

%%K t07-c1-12-1 p
12. --- Quella idea dal santeri ha fatg s'algordar al chalzaner da las persunas che
mes bab aveva enconuschì e ch'èn mortas l'onn passà. «Quantas \VUgras{fossas}
\textsuperscript{(21)}, mes bun signur, èn vegnidas ad ir a las autras, sut la
\VUgras{cipressa} \textsuperscript{(22)}! La mort ha seà noss vegl notar. Tut ils
vitgadins han assistì a sia \VUgras{sepultura}.»

%%K t07-c1-13-1 p
13. --- «Qua è ses successur, che patinescha sin l'auatg. El è gist vegnì a laschar
reparar da mai la curriaglia da ses \VUgras{patin} \textsuperscript{(26)}, ch'era
rutta.»

%%K t07-c1-14-1 p
14. --- «Qua è era, sin la riva da l'auatg, la nova \VUgras{guardia champestra}
\textsuperscript{(27)}. El è in ex-schandarm che spaventa ils mats naschas dal vitg.
Els fugian immediatamain cur ch'els vesan sias \VUgras{gambieras}
\textsuperscript{(29)} e ses \VUgras{kepi} \textsuperscript{(28)}.»

%%K t07-c1-15-1 p
15. --- «Ha! qua è il vegl Josef che maina ina \VUgras{charrettada da fasch da lain}
\textsuperscript{(30)} per il pastrizer. --- Quai è ver, ha ditg mes bab, jau
enconusch ses giuf da \VUgras{muls} \textsuperscript{(31)}. Jau hai propi da discurrer
cun el, jau vegn a profitar da l'occasiun. Sin renveser, bab Jachen.»

%%K t07-c1-16-1 p
16. --- Gist cur ch'nus essan sortids, han ils mats dal vitg bandunà la scola e sa
prescheschà currend a la \VUgras{sbrischanera} \textsuperscript{(33)}, cun il ristg da
crudar e da sa rumper in member en la \VUgras{crudada} \textsuperscript{(34)}. In d'els
ha prendì sia \VUgras{sliusla} \textsuperscript{(32)} tar il charrader, ha mess en ella
in da ses cumpogns ed è partì currend.

%%K t07-c1-17-1 p
17. --- Auters èn sa prescheschads ad in \VUgras{mantun da naiv} \textsuperscript{(37)}
sper la \VUgras{punt} \textsuperscript{(40)} ed han cumenzà a bombardar l'\VUgras{um da
naiv} \textsuperscript{(39)} ch'els avevan fatg il di avant ed al qual els han mess in
chapè, ina pippa ed in satg da scola en bandulera.

%%K t07-c1-18-1 p
18. --- Els sa fan pauc quitads dal vent che glatscha e dal fraid. Ils pli blers
n'han betg in \VUgras{schal} \textsuperscript{(35)} e, per sa stgaudar suenter la
battaglia cun \VUgras{balluccas da naiv} \textsuperscript{(38)}, als basta in maun
plain \VUgras{chastognas} \textsuperscript{(42)} fitg chaudas, cumpradas da la veglia
\VUgras{vendidra} Jachina, che trembla dal fraid sut ses \VUgras{schalin}
\textsuperscript{(43)} memia sutgl.

%%K t07-tit-5 sub
\VUsaut{3.0mm}
\VUcentre{12.0pt}{0}{\textit{Segunda scena.}}
\VUsaut{3.0mm}

%%K t07-tit-6 sub
\VUpk{10.2pt}{40}{Il bain rustical durant la primavaira.}
\VUsaut{4.0mm}

%%K t07-c2-01-1 p
1. --- Malgrà tut ils plaschairs da l'enviern salidain nus cun profunda legria il
return da la \VUgras{primavaira}.

%%K t07-c2-01-2 p
Tge maletg allegrant che dat ina curt da bain, la saira, en il mais da matg!

%%K t07-c2-02-1 p
2. --- A l'orizont sa metta il \VUgras{sulegl} \textsuperscript{(1)} plaunet, inundond
la \VUgras{champagna} \textsuperscript{(3)} cun ses \VUgras{radis} \textsuperscript{(2)}
purpurins. Il diligent \VUgras{aradur} \textsuperscript{(6)} sa pressa d'arar ses
\VUgras{chomp} \textsuperscript{(4)}.

%%K t07-c2-02-2 p
Cun ses \VUgras{arader} \textsuperscript{(5)} agiuntà ad in vigurus \VUgras{chaval}
\textsuperscript{(7)} tira el il \VUgras{surtg} \textsuperscript{(8)}, en il qual il
\VUgras{semnader} \textsuperscript{(9)} vegn a bittar cun il maun plain il
\VUgras{sem} che vegn a dar las spias d'aur.

%%K t07-c2-03-1 p
3. --- Ils \VUgras{sieaders} \textsuperscript{(11)} pervis cun lur faultschs han seà
l'erva dal prà, ils fanaders han sitgentà il \VUgras{fain} \textsuperscript{(10)} al
volvend cun \VUgras{furtgas} \textsuperscript{(12)} e rastels, e qua returnan els.

%%K t07-c2-03-2 p
Ins van avant il pesant char tratg da bovs; els portan lur utensils sin la spatla.
Auters, pli malpigns, sa èn installads confortablamain sin il fain che dat bun
odur.

%%K t07-c2-04-1 p
4. --- Passond fan els in salid amiaivel al \VUgras{giardinier} \textsuperscript{(16)}
occupà en l'\VUgras{iert da fritga} \textsuperscript{(13)} a taglier ils \VUgras{arbers
da fritga} ed a nuglier ils \VUgras{pirers} \textsuperscript{(14)}. Ils
\VUgras{brognaclers} \textsuperscript{(15)} cumenzan a flurir e, en l'\VUgras{iert da
verduras} \textsuperscript{(17)} bain letomà, van las verduras, ils cabis e las
salatas gia bain.

%%K t07-c2-04-2 p
Sin il mirin sblatta in bel \VUgras{pavun} \textsuperscript{(18)} sia \VUgras{cua}, per
laschar admirar sias \VUgras{plimas} bellissimas.

%%K t07-tit-7 sub
\VUpk{10.2pt}{40}{La curt dal bain.}
\VUsaut{3.0mm}

%%K t07-c2-05-1 p
5. --- Las portas da la \VUgras{stalla dals chavals} \textsuperscript{(19)} èn avertas
ed ins vesa il rastè ed il \VUgras{letg da palia} \textsuperscript{(23)} da la
\VUgras{juvla} \textsuperscript{(21)}, ch'il \VUgras{servent dals chavals}
\textsuperscript{(20)} ha gist disgiuntà. Il \VUgras{poltrin} \textsuperscript{(22)},
uschè comic cun sias lungas chommas, suonda sia mamma sagliond.

%%K t07-c2-06-1 p
6. --- Sper il \VUgras{puzz} \textsuperscript{(29)} van las \VUgras{vatgas}
\textsuperscript{(25)} a baiver en il \VUgras{trog} \textsuperscript{(26)} da lain che
las serva sco bavraria. Il \VUgras{vadè} \textsuperscript{(24)} profitescha da
l'occasiun per tetgar, ed il \VUgras{tor} \textsuperscript{(27)}, che senta il bun odur
dal fain, bragia allegramain ed auza cun fiereza il chau cun ses \VUgras{corns}
\textsuperscript{(28)} pussants.

%%K t07-c2-07-1 p
7. --- Tuts lavuran. La \VUgras{serventa} \textsuperscript{(32)} tegna en maun la
\VUgras{corda} \textsuperscript{(31)} dal puzz. Ella trai aua cun ina \VUgras{sadella}
\textsuperscript{(30)} per abeverar il muvel. Sper il \VUgras{granari}
\textsuperscript{(33)} rastella in um la palia, ed avant la porta da la
\VUgras{chasa} \textsuperscript{(34)} fila ina veglia \VUgras{filadra}
\textsuperscript{(35)}, cun sia roda da filar e sia \VUgras{conocla}, il
\VUgras{glin} u il \VUgras{chanv} sco en il temp vegl.

%%K t07-tit-8 sub
\VUsaut{3.0mm}
\VUpk{10.2pt}{40}{Il pur.}
\VUsaut{3.0mm}

%%K t07-c2-08-1 p
8. --- Il \VUgras{pur} \textsuperscript{(36)} dirigia tut quellas lavurs e dat
instrucziuns a ses servents. El recumonda da metter sut tetg l'\VUgras{erpi}
\textsuperscript{(40)} ed il \VUgras{sap} \textsuperscript{(39)}, ch'ins ha emblidà
sper la \VUgras{chamona} \textsuperscript{(38)} dal \VUgras{chaun}
\textsuperscript{(37)}.

%%K t07-c2-09-1 p
9. --- Ses clums spaventan ina \VUgras{columba} \textsuperscript{(41)}, che sgola cun
tut sias alas en la \VUgras{columbera} \textsuperscript{(42)}, e las \VUgras{giaglinas}
\textsuperscript{(43)}, che sa pressan da turnar en la \VUgras{giaglinera}
\textsuperscript{(44)}.

%%K t07-c2-10-1 p
10. --- Avant la \VUgras{nursera} \textsuperscript{(48)}, qua è il vegl
\VUgras{pastur} \textsuperscript{(45)}, che maina enavos sia \VUgras{scossa}
\textsuperscript{(49)}. Tegnend ses \VUgras{bastun} \textsuperscript{(46)} da pastur,
ch'al manca mai, sco era sia \VUgras{blusa} \textsuperscript{(47)} sblatgida, maina el
en la stalla \VUgras{muntuns} \textsuperscript{(50)}, \VUgras{nursas}
\textsuperscript{(53)}, \VUgras{agnels} \textsuperscript{(54)}, \VUgras{chauras}
\textsuperscript{(51)} e \VUgras{chavrins} \textsuperscript{(52)}. Ses fidel
\VUgras{chaun} \textsuperscript{(55)} Argus vegn l'ultim e sfrunta cun ses blers la
diligenza da quels che restan enavos.

%%K t07-c2-11-1 p
11. --- Tut quel canaster e quel muvament na disturban betg la quietezza da la buna
\VUgras{vatga da latg} \textsuperscript{(56)}, che fa tintinar sia \VUgras{plumpa}
\textsuperscript{(57)}, entant ch'ins la mulscha avant la porta da la \VUgras{stalla}
\textsuperscript{(58)}.

%%K t07-c2-12-1 p
12. Ella para da chapir ch'ella sto dar ses latg, per che la \VUgras{pura}
\textsuperscript{(60)} possia preparar \VUgras{paintg} en la \VUgras{zangola}
\textsuperscript{(61)} u ils excellents \VUgras{chaschiels} \textsuperscript{(64)}, che
sguttan da l'assa messa sin la \VUgras{tina} \textsuperscript{(63)}.

%%K t07-c2-13-1 p
13. --- Tut la \VUgras{lataria} \textsuperscript{(59)} vegn tegnida fitg netta dal
\VUgras{servent dal bain} \textsuperscript{(65)}, ch'ha gist lavà ils \VUgras{vaschs da
latg} \textsuperscript{(62)} en la gronda sadella ch'el porta en maun.

%%K t07-tit-9 sub
\VUsaut{3.0mm}
\VUpk{10.2pt}{40}{La pulaglia.}
\VUsaut{3.0mm}

%%K t07-c2-14-1 p
14. --- Cun ses gross zoccels da lain chamina el in pau plimpamain, e fa uschia fugir
il \VUgras{tatschin} \textsuperscript{(68)}, las \VUgras{andas} \textsuperscript{(66)},
ils \VUgras{ands} \textsuperscript{(67)} e las \VUgras{auchas} \textsuperscript{(69)},
che avran largiamain \VUgras{lur pichel} \textsuperscript{(70)} e criden malquiets.

%%K t07-c2-14-2 p
Il \VUgras{gial} \textsuperscript{(71)}, pli battaglius, auza sia cotschna
\VUgras{cresta} \textsuperscript{(72)}, sbatta las plimas da sia \VUgras{cua}
\textsuperscript{(73)} e fa udir in «kikeriki» sunant.

%%K t07-c2-15-1 p
15. --- La \VUgras{giaglina cluccia} \textsuperscript{(74)} tschertga sin il
\VUgras{letom} \textsuperscript{(81)} cun ses \VUgras{pulschins}
\textsuperscript{(75)}. Il \VUgras{purtgrin} \textsuperscript{(78)} fugia tar sia
mamma, l'enorma \VUgras{scrofa} \textsuperscript{(77)} che nus vesain sper il sti dals
portgs.

%%K t07-c2-16-1 p
16. --- En lur cumpartiments maglian ils \VUgras{cunigls} \textsuperscript{(79)}
avidamain l'\VUgras{erva} \textsuperscript{(80)} delicata che la figlia dal pur als
porta. Gionina ha fitg gugent ses cunigls e, mintga di, suenter avair prendì ils
frais \VUgras{ovs} \textsuperscript{(76)} en la giaglinera, vegn ella a visitar ses
pitschens amis.

\VUsaut{6.0mm}
\VUfilet{22mm}
